\section{Invertibility and Non-Erasure}

\subsection{Finite-stage invertibility}

All eigenvalues of \(Q\) are strictly positive:
\[
98,\quad18,\quad3,\quad2.
\]
Therefore \(Q\) and \(\widehat{Q}\) are invertible.

Their determinants are
\[
\det Q
=
98\cdot18^5\cdot3^4\cdot2^5
>
0
\]
and
\[
\det\widehat{Q}
=
\left(\frac{9}{49}\right)^5
\left(\frac{3}{98}\right)^4
\left(\frac{1}{49}\right)^5
>
0.
\]

\begin{proposition}[Finite-stage non-erasure]
For every finite integer \(k\ge0\) and every
\(x,y\in\mathbb R^{15}\),
\[
\widehat{Q}^{\,k}x
=
\widehat{Q}^{\,k}y
\quad\Longleftrightarrow\quad
x=y.
\]
\end{proposition}

\begin{proof}
Every finite power of an invertible matrix is invertible.
\end{proof}

Thus no two complete linear sector states are merged by a finite
number of normalized quadratic aggregation steps.

\subsection{An invertible sequence with a noninvertible limit}

The contrast-concentration theorem gives
\[
\widehat{Q}^{\,k}
\longrightarrow
\Pzero.
\]
The limiting operator \(\Pzero\) has rank one and is not invertible.

Therefore Paper 14 exhibits the following exact finite-dimensional
phenomenon:
\[
\begin{array}{c}
\text{every finite iterate is invertible},\\[0.3em]
\text{the operator limit is many-to-one}.
\end{array}
\]

There is no contradiction.  The inverses
\[
\widehat{Q}^{-k}
\]
become increasingly ill-conditioned on centered directions as
\(k\) grows.

The finite state remains mathematically recoverable when the exact
iterate and exact output are retained.  Recovery becomes
progressively sensitive to any loss of centered precision.

\subsection{Exact distinction and effective resolution}

Let
\[
d=x-y\neq0
\]
be a complete distinction vector.  Then
\[
\widehat{Q}^{\,k}d\neq0
\]
for every finite \(k\).

If \(d\) is centered, however,
\[
\left\|
\widehat{Q}^{\,k}d
\right\|_2
\le
\left(\frac{9}{49}\right)^k
\|d\|_2.
\]

Thus a complete distinction may remain exact while its centered
registered amplitude becomes arbitrarily small.

For a declared resolution threshold
\(\varepsilon>0\), there may exist a finite \(k\) for which
\[
0
<
\left\|
\widehat{Q}^{\,k}d
\right\|_2
<
\varepsilon.
\]
The distinction is still present in exact arithmetic but is no
longer resolved by that thresholded readout.

This motivates the phrase
\emph{effective registration degeneracy}.  It must not be confused
with exact equality of complete states.

\subsection{Scalar readout loss}

The centered scalar
\[
R_{15}(x)
=
\|\Pperp x\|_2
\]
is itself many-to-one.  Distinct centered vectors may share the same
norm.  Therefore even before any iteration,
\[
R_{15}(x)=R_{15}(y)
\]
does not imply
\[
x=y.
\]

The normalized contrast fraction
\[
\zetaRel(x)
=
\frac{R_{15}(x)^2}{\|x\|_2^2}
\]
forgets still more structure: common scale, centered direction, and
all lifted data absent from the fifteen-component register.

Consequently, three levels must remain separate:

\begin{enumerate}[label=(\roman*)]
\item exact complete sector state;
\item fifteen-component incidence register;
\item scalar contrast readout.
\end{enumerate}

Loss at level (iii) is not erasure at level (i).

\subsection{Relation to retained history}

The earlier finite mechanics supplies stronger distinctions than the
quadratic sector register alone.  Covering holonomy, deck receipts,
integer refinements, and registered-history winding can distinguish
states or histories that have already become identical under a
coarser projection.

The present quadratic theorem does not yet act on those richer
objects.  It supplies the contrast-concentrating side of the proposed
saturation architecture.  The retained-history theorem supplies
examples of complete distinction surviving projection.

A full internal saturation theorem would require one admitted native
construction in which both occur in the same declared diagram:
\[
\begin{array}{ccc}
\text{complete admitted state}
&\longrightarrow&
\text{retained receipt or winding}\\
\downarrow
&&
\downarrow\\
\text{contrast-concentrating register}
&\longrightarrow&
\text{coarse scalar or visible state}.
\end{array}
\]

The current theorem chain has not yet supplied that crosswalk.

\subsection{The exact conclusion}

The exact conclusion of Sections 4--7 is:

\begin{quote}
The normalized Thalean quadratic contracts every centered sector
direction by a factor at most \(9/49\) per iterate, converges to the
common-mode projection, and remains invertible at every finite
iterate.
\end{quote}

This is a theorem of finite linear algebra.

Its saturation interpretation is internal and conditional.

Its physical interpretation remains open.
